% ---------------------------------------------------------------------------
% Developer Note — A7.2 Toolchain Architecture & Runtime Integration
% Author: Borealis S. Hedling
% Altitude: A7.2 (Non-Dual Layer)
% Purpose: Describe the system-agnostic toolchain required to implement
%          collapse-surface repair, outline dependencies, and explain how
%          A7.2 integrates with Lean formalization and runtime execution.
% ---------------------------------------------------------------------------

\documentclass[11pt]{article}
\usepackage{amsmath, amssymb}

\title{Developer Note — A7.2 Toolchain Architecture and Runtime Integration}
\author{Borealis S. Hedling}
\date{October 2026}

\begin{document}
\maketitle

\section*{Overview}

This Developer Note outlines the system-agnostic toolchain required to implement
the A7.2 Collapse-Surface Repair Protocol. It describes the four-layer
architecture, the dependencies for each layer, and how the resulting artifacts
integrate with Lean formalization and runtime execution. This note is required
before implementing the A7.2 tools.

% ---------------------------------------------------------------------------
\section{Toolchain Overview}

The A7.2 repair pipeline consists of four interoperable layers:

\begin{enumerate}
  \item Collapse-Surface Detector (Python / FEniCS / Dedalus)
  \item Non-Dual Projection Engine (Rust / WASM)
  \item Collapse-Surface Serializer (JSON / Protocol Buffers)
  \item Dual-Compatible Lift Engine (Lean 4 / Mathlib)
\end{enumerate}

Each layer corresponds to a distinct altitude boundary and enforces specific
invariants.

% ---------------------------------------------------------------------------
\section{Layer 1: Collapse-Surface Detector}

\textbf{Purpose:} Identify broken geometry at A7 altitude.

\textbf{Dependencies:}
\begin{itemize}
  \item Python 3.11+
  \item FEniCS / Dolfin (FEM)
  \item Dedalus (spectral methods)
  \item PETSc / SLEPc (eigenvalue solvers)
\end{itemize}

\textbf{Detections:}
\begin{itemize}
  \item curvature singularities,
  \item spectral tears,
  \item modal discontinuities,
  \item holonomy fractures.
\end{itemize}

This layer produces raw collapse-surface candidates $\Sigma_{\mathrm{raw}}$.

% ---------------------------------------------------------------------------
\section{Layer 2: Non-Dual Projection Engine}

\textbf{Purpose:} Apply non-dual projection operators to reconcile broken geometry.

\textbf{Dependencies:}
\begin{itemize}
  \item Rust 1.80+
  \item WASM target (optional)
  \item Serde / Serde JSON
\end{itemize}

\textbf{Functions:}
\begin{itemize}
  \item non-idempotent projection,
  \item spectral tear closure,
  \item modal adjacency restoration,
  \item collapse-surface mapping.
\end{itemize}

This layer produces stable collapse surfaces $\Sigma$.

% ---------------------------------------------------------------------------
\section{Layer 3: Collapse-Surface Serializer}

\textbf{Purpose:} Encode collapse surfaces for reintegration.

\textbf{Dependencies:}
\begin{itemize}
  \item JSON Schema (Draft 2020-12)
  \item Protocol Buffers (optional)
\end{itemize}

\textbf{Serialized Fields:}
\begin{itemize}
  \item collapse coordinates,
  \item spectral tear metadata,
  \item modal adjacency restoration metadata,
  \item provenance reconstruction metadata.
\end{itemize}

This layer produces $\Sigma_{\mathrm{encoded}}$.

% ---------------------------------------------------------------------------
\section{Layer 4: Dual-Compatible Lift Engine}

\textbf{Purpose:} Reintegration of collapse surfaces into A6 dual algebra.

\textbf{Dependencies:}
\begin{itemize}
  \item Lean 4
  \item Mathlib 4
\end{itemize}

\textbf{Functions:}
\begin{itemize}
  \item dual-compatible lift,
  \item continuity restoration,
  \item adjacency restoration,
  \item kernel reconstruction,
  \item spectral stability enforcement.
\end{itemize}

This layer produces $\iota(\Sigma)$, a dual-compatible reintegration.

% ---------------------------------------------------------------------------
\section{Altitude Discipline}

The toolchain enforces altitude boundaries:

\begin{itemize}
  \item A7.0–A7.2: non-dual algebra, collapse surfaces, projection.
  \item A6.5: firewall preventing non-dual leakage.
  \item A6: dual algebra reintegration.
\end{itemize}

No non-dual operator may descend into A6 without passing through the lift engine.

% ---------------------------------------------------------------------------
\section{Lean Formalization Integration}

Lean formalization is used to:

\begin{itemize}
  \item verify dual-compatible lifts,
  \item enforce kernel reconstruction,
  \item prove continuity restoration,
  \item validate adjacency preservation,
  \item guarantee spectral stability.
\end{itemize}

Lean does not execute non-dual operators; it verifies the reintegration.

% ---------------------------------------------------------------------------
\section{Runtime Execution Integration}

Runtime execution follows:

\begin{enumerate}
  \item Python detects broken geometry.
  \item Rust reconciles geometry via non-dual projection.
  \item JSON serializes collapse surfaces.
  \item Lean verifies reintegration.
  \item A6 runtime resumes.
\end{enumerate}

This ensures drift-neutral, altitude-safe execution.

% ---------------------------------------------------------------------------
\section{Conclusion}

This Developer Note defines the toolchain architecture, dependencies, and
altitude boundaries required to implement A7.2 collapse-surface repair. It also
explains how Lean formalization and runtime execution integrate with the
non-dual repair pipeline. This note must be completed before building the tools.

\end{document}
